\section{Como se distribuem os indicadores}
\label{sec:univariada}

A Tabela~\ref{tab:univariada} resume as distribuições dos sete
indicadores e as Figuras~\ref{fig:univA} a~\ref{fig:univC} trazem os
histogramas. Antes
dos números, um esclarecimento que simplifica todo o restante do
documento: a \textbf{mortalidade ajustada é, na prática, idêntica à
mortalidade geral} neste painel (a exclusão dos óbitos fetais, maternos
e neonatais quase não altera as taxas; a concordância entre as duas é
total, com correlação de 1,000). Por isso ela aparece na tabela, mas as
figuras mostram apenas a mortalidade geral, e tudo o que for dito vale
igualmente para a versão ajustada.

Em palavras simples, o que os dados mostram: o hospital típico do
painel tem \textbf{mortalidade em torno de 5,4\%} (metade dos hospitais
está entre 3,0\% e 7,0\%), \textbf{permanência média de 4,4 dias} e
\textbf{custo de cerca de R\$~1.071 por internação encerrada}. Mas os
extremos importam. Na mortalidade há um segundo grupo de hospitais com
taxas próximas de zero (maternidades e especializados) e \textbf{1,5\%
das observações têm mortalidade exatamente zero}. No tempo de
permanência existe uma \textbf{subpopulação que se acumula num teto de
30,5 dias}: são hospitais de perfil psiquiátrico ou de cuidados
prolongados, em que a permanência longa é característica do serviço (o
teto reflete o ciclo mensal de faturamento da AIH), e não erro de
registro. No custo, a cauda direita é longa: enquanto a mediana é
R\$~1.071, o máximo chega a R\$~27.888, e é por isso que a média
(R\$~1.475) fica 38\% acima da mediana. A fração de alta complexidade é
o indicador mais concentrado: \textbf{34,4\% dos hospitais ano não
realizam nenhuma internação de alta complexidade}. Na ocupação,
\textbf{14,7\% das observações de internação ultrapassam 100\%}
(hospital operando acima da capacidade cadastrada, informação real de
sobrecarga) e a UTI combina \textbf{25,5\% de zeros} (hospitais sem UTI
ativa) com um máximo impossível de 875\%, examinado na
Seção~\ref{sec:qualidade}.

\begin{table}[H]
\centering
\caption{Resumo estatístico dos sete indicadores (3.454 observações de
hospital e ano; mortalidades e alta complexidade em proporção;
ocupações em percentual). DP é o desvio padrão, medida do espalhamento
dos valores; p90 é o valor abaixo do qual estão 90\% das observações;
assimetria positiva indica cauda comprida à direita.}
\label{tab:univariada}
\small
\begin{tabular}{lrrrrrrr}
\toprule
Indicador & Média & DP & Mediana & p90 & Máximo & Assimetria & Zeros \\
\midrule
Mortalidade geral      & 0,0517 & 0,0319 & 0,0536 & 0,0878 & 0,2602 & 0,72 & 1,5\% \\
Mortalidade ajustada   & 0,0517 & 0,0319 & 0,0535 & 0,0878 & 0,2602 & 0,72 & 1,6\% \\
TMP (dias)             & 5,99   & 5,75   & 4,42   & 7,98   & 30,50  & 3,33 & 0,03\% \\
Custo por saída (R\$)  & 1.475,00 & 1.673,59 & 1.071,27 & 2.459,93 & 27.887,60 & 6,59 & 0\% \\
Fração alta complex.   & 0,0731 & 0,1530 & 0,0020 & 0,2215 & 1,0000 & 3,41 & 34,4\% \\
Ocupação internação (\%) & 71,42 & 27,72 & 73,46 & 103,83 & 236,61 & 0,05 & 0\% \\
Ocupação UTI (\%)      & 56,83  & 45,29  & 67,15  & 96,10  & 875,21 & 2,72 & 25,5\% \\
\bottomrule
\end{tabular}
\end{table}

\begin{figure}[H]
\centering
\subfig{fig_ae_01_hist_mort_all.png}{Mortalidade geral}
\subfig{fig_ae_01_hist_tmp.png}{TMP}
\caption{Distribuição da mortalidade geral e do tempo médio de
permanência (314 hospitais, 2015 a 2025). Em cada gráfico, a linha
tracejada vertical é a mediana, a pontilhada é a média e os valores de
pico e máximo estão anotados.}
\label{fig:univA}
\end{figure}

\comoler{cada barra mostra quantas observações caem naquele valor do
indicador. Na mortalidade, note as duas corcovas: a da esquerda, perto
de zero, são as maternidades e especializados; a principal fica perto
de 5,5\%. No TMP, quase todos os hospitais estão entre 3 e 8 dias, mas
há um pequeno grupo isolado junto ao máximo de 30,5 dias, os hospitais
de longa permanência.}

\begin{figure}[H]
\centering
\subfig{fig_ae_01_hist_custo_saida.png}{Custo por saída}
\subfig{fig_ae_01_hist_pct_alta_complex.png}{Fração alta complexidade}
\caption{Distribuição do custo por saída e da fração de alta
complexidade.}
\label{fig:univB}
\end{figure}

\comoler{no custo, a concentração fica perto de R\$~1.000 e a cauda se
estica até quase R\$~28.000, o que explica a média (pontilhada) ficar
bem à direita da mediana (tracejada). Na alta complexidade, a barra
gigante em zero mostra os 34,4\% de hospitais ano sem nenhuma
internação desse tipo.}

\begin{figure}[H]
\centering
\subfig{fig_ae_01_hist_ocupacao_internacao.png}{Ocupação de internação}
\subfig{fig_ae_01_hist_ocupacao_uti.png}{Ocupação de UTI}
\caption{Distribuição das taxas de ocupação de internação e de UTI, em
escala percentual.}
\label{fig:univC}
\end{figure}

\comoler{a linha pontilhada vertical marca 100\%: tudo à direita dela é
operação acima da capacidade cadastrada. Na internação, isso ocorre em
14,7\% das observações. Na UTI, a barra em zero são os hospitais sem
UTI ativa, e a cauda que se arrasta até 875\% é o problema de cadastro
examinado na Seção~\ref{sec:qualidade}.}
